\documentclass[11pt]{article}
\usepackage[margin=1in]{geometry}
\usepackage{amsmath,amssymb,hyperref}
\title{Goal revision with an outstanding right to use a lane}
\date{24 November 2026}
\begin{document}
\maketitle

At tick 2 a warehouse storage operator changes the same goal in two executions: return the load from B to A by tick 3 and keep it there. In both executions the load is on B and A is empty. The return succeeds in one tick only in the execution without the outstanding right. In the other, an independent crew can lawfully start its already promised use of A at tick 2 and retain that lane until tick 4. Even a scheduler with perfect information cannot guarantee both the original promise and permanent restoration by tick 3. The same installation guarantees both with a tick-5 restoration deadline. Lowering the promise's priority in the controller without authority is not a repair.

\section{Authority, physical state, and the reached cuts}
The storage operator orders the initial move from A to B at tick $-2$; it completes at $-1$. In the right-bearing run the facility grants and the service crew accepts an irrevocable one-use clearance on A at $-1$. This is a binding promise of admission at any lawful start, not a message forecasting a grant. The independent crew chooses whether and when to exercise; its one use must observe the allowed start and dwell relation. No use in the common comparison has started before the revision. In the control run the allocation is identical except that no clearance is granted. Thus the model reaches the two revision cuts from the same stock before the initial move; it does not postulate two terminal yard photographs.

At effective tick $R$, a verified decision by the storage operator replaces its own destination goal and seals further issuance of this lane's service rights. An unauthenticated revision does neither. The priority rule at a tick is due physical releases, closure of unexercised rights after their last eligible start, then revision, authenticated local notice and the controller's move decision, then an independently scheduled old start. A start at the exact exclusive expiry $E$ is invalid. If admitted by the clearance relation, a lawful old start at $E-1$ remains possible in that phase, even after a revision or the delivery of its notice. Simultaneous release and a new start are allowed because the old occupancy interval is half open. No receiver initiates the return before its own authenticated notice at $N\ge R$, although the old goal's change and sealing take effect at $R$. The actor has the local clock and its load state. It receives the maximum admissible timing envelope, and in the primary policy a trusted sealed-roster bit in the revision notice. Hidden authorization of the independent crew is not an admissible local observation without that receipt. A receipt scoped to the current sealed revision may also certify permanent freedom: every old right is irrevocably spent, waived or expired, and any associated physical occupation has actually ended.

The operator can withdraw the old destination obligation but cannot waive the crew's independent clearance, shorten an active occupation, or cancel a lawful old start. Only the holder may explicitly waive its unexercised clearance; a waiver does not settle active physical work. Rights are not transferable in this envelope. The uniform policy instead lets old work start, finish and physically release or lets its unexercised authorization expire. At expiry a started use may still occupy A: with $E=3$ the use on $[2,4)$ persists at tick 3. The new return carries a fresh physical reservation from its start until placement and leaves A occupied thereafter. Settlement of an old identity is not a statement that the lane will stay empty after the new reservation. There is no substitute lane, temporary return or virtual restoration; the moving load and crew require exclusive access to A. The facility guarantees bounded duration and reliable clock ticks. No communication, report or authority fault beyond the declared lateness or missing roster evidence is assumed.

\section{Exact window and locally executable transition}
Let $P$ be the nonempty finite set of jointly admissible start/dwell pairs $(a,d)$, with $0\le a<E$ and $d\ge1$. Here a right known to have been granted and with no closure evidence may have started before or after the goal revision. A complete run can also omit exercise. Let $C\ge1$ be the return duration. With a valid roster bit in the notice and no earlier discharge certificate, the tight worst completion is
\[
 T_0=N+C \quad\text{if no right exists},\qquad
 T_1=\max\{N,\max_{(a,d)\in P}(a+d)\}+C \quad\text{otherwise}.
\]
Before $N$ there is no local authority to start. Choose a pair attaining the inner maximum. Starting restoration before that occupation ends either collides with it, or, if restoration precedes the independent start, makes the permitted use impossible since A must stay occupied. A physical interlock that refuses the old use trades the conflict for a breach of the original service promise. These are necessary constraints even for a full-state scheduler; no particular tie rule can remove them while keeping faithful service. The installed controller is to wait for the signed notice, then start immediately if the roster is empty; otherwise start at the computed upper bound, or earlier on a trusted permanent-clearance receipt. There is no pending occupation or possible future exercise at the bound. If the maximum release occurs before $E$, future exercise is already ruled out by the declared joint start/dwell relation: since $d\ge1$, any future eligible start would imply a release beyond the maximum. The guard closes a still-unexercised token immediately after its last eligible start tick, not at a later administrative expiry that adds no eligible start. The replay includes a sparse-window control for this distinction. On receipt of exact closure the same argument applies even before the bound; the wait is conservative where no such proof is available. Exclusivity preserves the physical model through new placement, and a settled old job remains settled.

For the start-by-expiry rule $P=\{0,\ldots,E-1\}\times\{1,\ldots,D\}$ the bound with a live right is $\max(N,E-1+D)+C$. A pair attaining the two separate maxima is load-bearing, not a statistical independence assumption. For the finish-by-expiry rule restrict the same set to $a+d\le E$; this gives $\max(N,E)+C$. At $E=3,D=2,C=1,N=2$, both allow last start 2 and duration 2, but finish-by permits these only on different histories; restoration guarantees differ (5 versus 4). Changing positive control's deadline is explicit, not proof that the tick-3 request has been met.

The vacancy-only control attempts at tick 2, restores at tick 3 and refuses the intact old right at tick 2. This is recorded as an unsafe service policy even though no overlap is permitted. If the roster bit is hidden, the no-right and outstanding-right executions can be locally identical at tick 2. A policy cannot guarantee the tick-3 return on the no-right execution while preserving the binding promise in the other: the necessary tick-2 action is common to both. Waiting according to the admitted right envelope protects both but misses the no-right opportunity. A timely authenticated permanent-clearance receipt separates those histories. The executable check also distinguishes permission to start following a revision from old clearance settlement, an authorized pre-start waiver from unilateral cancellation, and known no-right status with a late notice from immediate local authority.

\section{Evidence and limits}
\texttt{replay.py} models actual forward movement, issue and acceptance, notice, scheduled starts, releases and returns; an exact reference scans candidate start slots against every permitted old occupancy interval. The replay checks 2,160 parameter/envelope combinations and 9,045 reached sweep executions and checks the above controls. The no-right and unexercised runs share the same physical stock and reporting phase at the revision. The same implementation preserves the old service when it meets the new achievable deadline and reports a missed new deadline rather than silently refusing owed work in the impossible case. The direct baseline is a reservation-and-rights guard with the same timing and evidence. The comparison borrows \href{https://github.com/eternal-november/eos-research-programmes/pull/354}{the distinction between present stock and an interval assertion} and \href{https://github.com/eternal-november/eos-research-programmes/pull/352}{identity-bound documentary settlement}, not their theorems as black-box guarantees. \href{https://www.cs.cornell.edu/andru/cs711/2002fa/reading/sagas.pdf}{Sagas} give a prior model of semantic compensation instead of erasing visible effects; \href{https://web.stanford.edu/class/archive/cs/cs240/cs240.1076/readings/89-leases.pdf}{leases} give bounded exercise authority. Neither precedent lets this operator erase an already accepted physical duty merely by revising its own goal. No scientific novelty, proof-reuse saving, unbounded timing guarantee, physical validation or universal goal-revision result is claimed.
\end{document}
